% =====================================================================
%  SURFACE — Evaluation and Results
%  Structure modelled on YuraScanner (NDSS'25): an "Experimental Setup"
%  block (ground truth, testbed categories, baselines, modes, environment),
%  then one subsection per research question, each written as a
%  "Methodology -> Results" pair, exactly as YuraScanner does for RQ2–RQ5.
%
%  Required packages: booktabs, multirow, array.  (xcolor+\cellcolor only
%  for the highlighted winning row, as in the current Table 7.)
% =====================================================================

\section{Evaluation and Results}
\label{sec:eval}

We evaluate \textsc{Surface} against the four research questions of
Section~\ref{sec:problem}. The evaluation is designed to answer, concretely:
(i)~can a black-box scanner decide from a seed URL alone whether an
application exposes a reachable, model-backed AI surface, and with what
precision and recall (RQ1); (ii)~can it separate application-owned
conversational surfaces from third-party widgets and human-operated channels
so that active testing stays in scope (RQ2); (iii)~does discovery followed by
validation surface attackable AI interfaces and conventional vulnerabilities
that existing tools miss under an equal scanning budget (RQ3); and (iv)~can
prompt-injection findings meet the same evidence bar as conventional findings
while keeping false positives low (RQ4). Throughout, we contrast
\textsc{Surface} with both traditional web scanners and AI red-teaming tools,
and we report results in two complementary modes: an \emph{oracle-endpoint}
mode that isolates testing capability and a \emph{black-box discovery} mode
that measures end-to-end discovery from a URL.

% =====================================================================
\subsection{Experimental Setup}
\label{sec:eval:setup}

% ---------------------------------------------------------------------
\subsubsection{Ground-Truth Construction}
\label{sec:eval:gt}

A central requirement of our evaluation is that ground truth is established
\emph{independently} of any scanner output. For every target we build a label
record by manual inspection of the running application, its rendered
interface, its network traffic, and—where available—its source and container
image. Two authors annotated each target independently and reconciled
disagreements through discussion; a random cross-validation sample showed no
residual disagreement. The complete per-target records are released as an
artifact and summarised in the Appendix.

Each target is recorded under an anonymised identifier and classified as
either \emph{AI-enabled} or \emph{conventional}. A set of fields then captures
whether the target exposes a live AI surface and, if so, its type, whether
that surface is reachable without credentials, whether it is application-owned
or third-party, and whether it is a human-operated support channel rather than
a model. Two further fields flag the difference between genuine and merely
apparent AI exposure: whether the target only \emph{mentions} AI without a live
model, and whether the only evidence is an imported SDK or a marketing hint.
Finally, each record carries evaluation annotations: the known conventional
vulnerabilities, the known prompt-injection behaviours, and the expected false
positives for that target. Table~\ref{tab:gt-schema} defines the labelling
schema and the permitted values for each field.

\begin{table}[t]
\centering
\caption{Ground-truth labelling schema. Each target is annotated
independently of scanner output. ``annotation'' fields are lists used to score
precision and recall.}
\label{tab:gt-schema}
\small
\begin{tabular}{@{}lp{3.9cm}@{}}
\toprule
\textbf{Field} & \textbf{Permitted values} \\
\midrule
\texttt{target\_id}            & anonymised ID (T01--T40) \\
\texttt{class}                 & \{AI-enabled, conventional\} \\
\texttt{live\_ai\_surface}     & \{yes, no\} \\
\texttt{surface\_type}         & \{REST inference API, GraphQL, chatbot, WebSocket, streaming, none\} \\
\texttt{reachable\_no\_creds}  & \{yes, no, partial\} \\
\texttt{application\_owned}    & \{yes, no, unknown\} \\
\texttt{third\_party\_widget}  & \{yes, no\} \\
\texttt{human\_operated}       & \{yes, no\} \\
\texttt{mentions\_ai\_only}    & \{yes, no\} \\
\texttt{sdk\_or\_marketing\_only} & \{yes, no\} \\
\midrule
\texttt{known\_web\_vulns}     & list of (class, path, param) \emph{(annotation)} \\
\texttt{known\_pi\_behaviours} & list of injection behaviours \emph{(annotation)} \\
\texttt{expected\_false\_pos}  & list of must-not-report items \emph{(annotation)} \\
\bottomrule
\end{tabular}
\end{table}

The labelling schema is only meaningful if we fix, in advance, what counts as a
genuine attackable AI surface as opposed to a superficial indicator.
Table~\ref{tab:surface-def} states this operational definition. Superficial
indicators—marketing claims, imported SDK scripts, or OpenAPI/GraphQL
references to endpoints that are not reachable—are classified as non-attackable
structural hints and never, on their own, trigger AI-specific testing. Concrete
model behaviour observed from a live probe—an OpenAI-compatible completion
object, or a chat UI that returns backend-generated text—is classified as a
live or conversational AI surface. Third-party embedded widgets and
human-operated support forms are recorded as part of the attack surface but are
marked out-of-scope for active testing absent explicit authorisation.

\begin{table}[t]
\centering
\caption{Operational definition of an attackable AI surface. Only live,
application-owned surfaces are eligible for active AI-specific testing.}
\label{tab:surface-def}
\small
\begin{tabular}{@{}p{3.5cm}ll@{}}
\toprule
\textbf{Observed evidence} & \textbf{Classification} & \textbf{Testable} \\
\midrule
Marketing / brand claim                         & structural hint   & no \\
Imported AI SDK script                          & structural hint   & no \\
OpenAPI/GraphQL ref to unreachable endpoint     & structural hint   & no \\
Live OpenAI-compatible completion object        & live AI surface   & yes$^\dagger$ \\
Chat UI returning backend text                  & conversational    & yes$^\dagger$ \\
Third-party embedded AI widget                  & out-of-scope      & no$^\ddagger$ \\
Human-operated support form                     & out-of-scope      & no$^\ddagger$ \\
\bottomrule
\end{tabular}
\\[2pt]
{\footnotesize $^\dagger$ when the surface is application-owned.\quad
$^\ddagger$ unless explicit authorisation is provided.}
\end{table}

% ---------------------------------------------------------------------
\subsubsection{Testbed and Experiment Categories}
\label{sec:eval:testbed}

We assembled a testbed of 40 containerised web applications, organised into
four categories of ten targets each. The categories are not cosmetic: each
isolates a distinct property that the evaluation must exercise, so that a
single aggregate number cannot hide a weakness.
Table~\ref{tab:categories} lists the categories, their size, and the role each
plays; the full application inventory with versions appears in the Appendix
(Table~\ref{tab:inventory}). \emph{Vulnerable web applications} are
AI-negative and carry documented conventional vulnerabilities; they measure
conventional detection and, crucially, false-positive suppression on targets
that contain no live model. \emph{AI/LLM applications} expose live,
application-owned inference endpoints and are the primary targets for AI
surface discovery and prompt-injection assessment. \emph{Chatbot /
conversational applications} expose rendered conversational interfaces, some
application-owned and some backed by third-party services, and exercise the
separation of detection from authorisation. \emph{Conventional applications
with widgets} are AI-negative but deliberately seeded with third-party widgets,
imported SDK scripts, and marketing copy; they are the precision stress test
against superficial AI indicators.

\begin{table}[t]
\centering
\caption{Testbed composition. Four categories of ten targets each isolate a
distinct evaluation property.}
\label{tab:categories}
\small
\begin{tabular}{@{}lcp{3.4cm}@{}}
\toprule
\textbf{Category} & \textbf{\#} & \textbf{Why it is necessary} \\
\midrule
Vulnerable web apps        & 10 & conventional recall/precision; AI-negative, tests FP suppression \\
AI / LLM applications      & 10 & live, owned inference endpoints; AI discovery + prompt injection \\
Chatbot / conversational   & 10 & owned vs.\ third-party interfaces; detection-vs-authorisation \\
Conventional + widgets     & 10 & SDK/marketing/third-party decoys; precision against superficial AI \\
\midrule
\textbf{Total}             & \textbf{40} & \\
\bottomrule
\end{tabular}
\end{table}

% ---------------------------------------------------------------------
\subsubsection{Baselines}
\label{sec:eval:baselines}

We compare \textsc{Surface} against two families of tools. The first is
traditional web scanners: OWASP ZAP~\cite{zap}, Burp Suite~\cite{burp}, Black
Widow~\cite{blackwidow}, YuraScanner~\cite{yurascanner}, and
Nikto~\cite{nikto}. These tools discover and test conventional endpoints but
have no notion of an AI surface; they establish the conventional-vulnerability
baseline and let us quantify the AI attack surface they miss. The second family
is AI red-teaming tools: PyRIT~\cite{pyrit}, garak~\cite{garak},
Promptfoo~\cite{promptfoo}, and DeepTeam~\cite{deepteam}. These tools test
prompt injection and related attacks \emph{given} an inference endpoint, but
assume the endpoint, request schema, and credentials are already known; they
establish the AI-testing baseline and let us isolate the contribution of
discovery.

% ---------------------------------------------------------------------
\subsubsection{Evaluation Modes}
\label{sec:eval:modes}

Because the two baseline families make opposite assumptions about what is known
at the start of a scan, we evaluate in two modes.
\emph{Oracle-endpoint mode}: a human supplies the correct AI inference endpoint
(and, where required, the request schema and credentials) to every tool, and we
measure payload success. This mode isolates testing capability and gives the AI
red-teaming tools their most favourable setting.
\emph{Black-box discovery mode}: every tool receives only the seed URL, and we
measure whether the system locates the AI surface at all before any payload is
sent. This mode reflects the deployment setting \textsc{Surface} targets and is
the one in which discovery is the differentiator.

% ---------------------------------------------------------------------
\subsubsection{Environment and Scanning Budget}
\label{sec:eval:env}

All targets and tools run as Docker containers on an isolated virtual network;
the application state is reset between runs to avoid interference. The language
model is served locally so that target URLs and discovered content never leave
the evaluation infrastructure. To keep the comparison fair, authentication is
handled uniformly: every tool, including \textsc{Surface}, receives the same
credentials or session cookies and the same scripted login sequence for targets
that require one. Each scan is bounded by a fixed time budget of
\textsc{XXX}~hours, and all reported costs (wall-clock time and token usage) are
measured under this budget.

% =====================================================================
\subsection{Metrics}
\label{sec:eval:metrics}

We report, per target: (i)~the number of confirmed and potential findings;
(ii)~the precision of confirmed findings, established by manual review, as a
proxy for the false-positive rate; (iii)~SPA coverage, measured as the number
of routes and forms discovered and tested on client-rendered targets; and
(iv)~end-to-end scan duration and token cost. For AI surface discovery we
additionally report the endpoint discovery rate (fraction of live AI surfaces
located in black-box mode) and, for conventional and prompt-injection findings
alike, true/false positives and negatives scored against the independent ground
truth of Section~\ref{sec:eval:gt}. Precision is $\mathit{TP}/(\mathit{TP}+
\mathit{FP})$, recall $\mathit{TP}/(\mathit{TP}+\mathit{FN})$, and accuracy
$(\mathit{TP}+\mathit{TN})/(\mathit{TP}+\mathit{TN}+\mathit{FP}+\mathit{FN})$
over the fixed decision grid defined in the Appendix.

% =====================================================================
\subsection{RQ1: AI Surface Discovery}
\label{sec:eval:rq1}

\textbf{Methodology.} Using the AI/LLM and Conventional+widgets categories, we
run \textsc{Surface} in black-box discovery mode and score each target's
predicted label (live AI surface / none, and surface type) against the
ground-truth schema. The Conventional+widgets category is the decisive test:
its targets advertise AI through SDK imports, OpenAPI references, and marketing
copy but expose no live model, so any AI-positive verdict there is a false
positive. We report precision, recall, and the confusion matrix over the
AI/non-AI decision, and we break down false positives by the superficial
indicator that triggered them.

\textbf{Results.} Table~\ref{tab:rq1} reports discovery performance.
\textsc{Surface} achieves a precision of \textsc{XXX}\% and recall of
\textsc{XXX}\% on the AI/non-AI decision, while the conventional and
red-teaming baselines either report no AI verdict or are driven to
\textsc{XXX}\% false positives by SDK and marketing signals. [Narrate the
false-positive breakdown and the surface-type confusion here.]

\begin{table}[t]
\centering
\caption{RQ1 --- AI surface discovery (black-box mode). Decision is
AI-enabled vs.\ conventional; FP driven by SDK/marketing/soft-404 signals.}
\label{tab:rq1}
\small
\begin{tabular}{@{}lccccc@{}}
\toprule
\textbf{System} & \textbf{TP} & \textbf{FP} & \textbf{FN} &
\textbf{Prec.\ (\%)} & \textbf{Rec.\ (\%)} \\
\midrule
ZAP / Burp / Nikto   & -- & -- & -- & -- & -- \\
Black Widow          & -- & -- & -- & -- & -- \\
YuraScanner          & -- & -- & -- & -- & -- \\
garak / PyRIT        & -- & -- & -- & -- & -- \\
\textbf{\textsc{Surface}} & \textsc{XXX} & \textsc{XXX} & \textsc{XXX} & \textsc{XXX} & \textsc{XXX} \\
\bottomrule
\end{tabular}
\end{table}

% =====================================================================
\subsection{RQ2: Conversational Surface Ownership and Scope}
\label{sec:eval:rq2}

\textbf{Methodology.} On the Chatbot/conversational category we evaluate two
decisions separately: \emph{detection}, whether a conversational AI interface
exists, and \emph{authorisation}, whether that interface is application-owned
and therefore eligible for active injection testing. Each target is labelled in
ground truth as application-owned, third-party widget, or human-operated. We
score \textsc{Surface}'s detection against the first label and its
authorisation decision against the second, and we verify that no active
injection payload is sent to any target labelled third-party or human-operated.

\textbf{Results.} Table~\ref{tab:rq2} reports detection and authorisation
accuracy and confirms that active testing was confined to application-owned
surfaces. \textsc{Surface} detects \textsc{XXX}\% of conversational interfaces
and correctly withholds active testing from \textsc{XXX}/\textsc{XXX}
third-party or human-operated targets. [Narrate any misclassifications and
their cause here.]

\begin{table}[t]
\centering
\caption{RQ2 --- detection vs.\ authorisation on conversational targets.
Active testing must be withheld from third-party and human-operated surfaces.}
\label{tab:rq2}
\small
\begin{tabular}{@{}lccc@{}}
\toprule
\textbf{Decision} & \textbf{Correct} & \textbf{Total} & \textbf{Acc.\ (\%)} \\
\midrule
Conversational detection      & \textsc{XXX} & \textsc{XXX} & \textsc{XXX} \\
Ownership / authorisation     & \textsc{XXX} & \textsc{XXX} & \textsc{XXX} \\
Out-of-scope correctly skipped& \textsc{XXX} & \textsc{XXX} & \textsc{XXX} \\
\bottomrule
\end{tabular}
\end{table}

% =====================================================================
\subsection{RQ3: Attack Surface Coverage}
\label{sec:eval:rq3}

\textbf{Methodology.} We measure coverage in both modes across all four
categories under the fixed scanning budget. In oracle-endpoint mode the correct
AI endpoint is supplied to every tool and we measure prompt-injection payload
success; this is the setting most favourable to the AI red-teaming baselines. In
black-box discovery mode only the seed URL is supplied and we measure whether
each tool reaches the AI surface at all. For conventional vulnerabilities we
additionally compute, per tool pair, the unique and shared surface
($A\setminus B$, $A\cap B$, $B\setminus A$) following prior
work~\cite{yurascanner,blackwidow}, so that orthogonal coverage is visible
rather than hidden in a single total.

\textbf{Results.} Table~\ref{tab:rq3} contrasts the two modes. In
oracle-endpoint mode the red-teaming tools achieve payload success comparable to
\textsc{Surface} (\textsc{XXX}\% vs.\ \textsc{XXX}\%), confirming that the
testing payloads are not the differentiator. In black-box discovery mode those
same tools reach \textsc{XXX}\% of AI surfaces because they cannot locate the
endpoint, whereas \textsc{Surface} reaches \textsc{XXX}\%. On conventional
targets \textsc{Surface} recovers the baseline coverage and additionally
confirms \textsc{XXX} AI-specific findings that no baseline reports under the
same budget. [Narrate the unique/shared coverage breakdown here.]

\begin{table*}[t]
\centering
\caption{RQ3 --- coverage in oracle-endpoint and black-box discovery modes.
Oracle mode isolates testing capability; discovery mode measures
end-to-end discovery from a URL.}
\label{tab:rq3}
\small
\begin{tabular}{@{}lcccccc@{}}
\toprule
& \multicolumn{3}{c}{\textbf{Oracle-endpoint mode}}
& \multicolumn{3}{c}{\textbf{Black-box discovery mode}} \\
\cmidrule(lr){2-4}\cmidrule(lr){5-7}
\textbf{System} & \textbf{Endpoints} & \textbf{PI success (\%)} & \textbf{Web vulns}
                & \textbf{Endpoints found} & \textbf{PI success (\%)} & \textbf{Web vulns} \\
\midrule
ZAP / Burp / Nikto        & n/a         & n/a         & \textsc{XXX} & \textsc{XXX} & n/a         & \textsc{XXX} \\
Black Widow / YuraScanner & n/a         & n/a         & \textsc{XXX} & \textsc{XXX} & n/a         & \textsc{XXX} \\
garak / PyRIT / Promptfoo / DeepTeam & given & \textsc{XXX} & n/a & \textsc{XXX} & \textsc{XXX} & n/a \\
\textbf{\textsc{Surface}} & given       & \textsc{XXX} & \textsc{XXX} & \textsc{XXX} & \textsc{XXX} & \textsc{XXX} \\
\bottomrule
\end{tabular}
\end{table*}

% =====================================================================
\subsection{RQ4: Evidence-Grounded Reporting}
\label{sec:eval:rq4}

\textbf{Methodology.} We evaluate whether prompt-injection findings meet the
same evidence bar as conventional findings and whether the differential
control--treatment protocol reduces false positives relative to single-response
confirmation. For each prompt-injection attempt we run the paired control and
treatment interactions and count the attack as successful only when the target
behaviour appears in the treatment but not the control. As an ablation we also
score the same attempts under single-response confirmation (treatment only) and
report the resulting false-positive rate. For conventional findings we reuse the
independent ground truth and report the confusion matrix per target.

\textbf{Results.} Table~\ref{tab:detection} reports conventional-vulnerability
detection against the benchmark targets, scored against ground truth.
\textsc{Surface} matches or exceeds the baselines on confirmed findings while
holding false positives low. The differential protocol reduces the
prompt-injection false-positive rate from \textsc{XXX}\% under single-response
confirmation to \textsc{XXX}\% [narrate the ablation here]; every reported
prompt-injection finding is backed by a reproducible control--treatment pair and
mapped to OWASP~LLM01:2025 and MITRE ATLAS.

\begin{table*}[t]
\centering
\caption{Vulnerability detection performance of YuraScanner, Black Widow, and
\textsc{Surface} on the benchmark applications. TP, FP, TN, and FN denote true
positives, false positives, true negatives, and false negatives. Precision is
$\mathit{TP}/(\mathit{TP}+\mathit{FP})$, recall $\mathit{TP}/(\mathit{TP}+
\mathit{FN})$, and accuracy $(\mathit{TP}+\mathit{TN})/(\mathit{TP}+
\mathit{TN}+\mathit{FP}+\mathit{FN})$.}
\label{tab:detection}
\small
\begin{tabular}{@{}llccccrrr@{}}
\toprule
\textbf{Scanner} & \textbf{Target} & \textbf{TP} & \textbf{FP} & \textbf{TN} &
\textbf{FN} & \textbf{Prec.\ (\%)} & \textbf{Rec.\ (\%)} & \textbf{Acc.\ (\%)} \\
\midrule
YuraScanner & DVWA & 2 & 0 & 1 & 15 & 100   & 11.76 & 16.67 \\
Black Widow & DVWA & 0 & 0 & 1 & 13 & NA    & 0     & 7.14  \\
\textsc{Surface} & DVWA & 2 & 0 & 0 & 0 & 100 & 100 & 100 \\
\midrule
YuraScanner & Hackazon & 0 & 0 & 1 & 7 & 0   & 0  & 12.50 \\
Black Widow & Hackazon & 0 & 0 & 7 & 0 & NA  & 0  & 0 \\
\textsc{Surface} & Hackazon & 2 & 3 & 0 & 3 & 40 & 40 & 25 \\
\midrule
YuraScanner & WebGoat & 0 & 0 & 1 & 11 & NA   & 0    & 8.33 \\
Black Widow & WebGoat & 0 & 0 & 1 & 11 & NA   & 0    & 8.33 \\
\textsc{Surface} & WebGoat & 1 & 2 & 0 & 11 & 33.33 & 8.33 & 7.14 \\
\midrule
YuraScanner & WrongSecrets & 0 & 0 & 6 & 1 & NA    & 0     & 85.71 \\
Black Widow & WrongSecrets & 0 & 0 & 6 & 1 & NA    & 0     & 85.71 \\
\textsc{Surface} & WrongSecrets & 2 & 1 & 5 & 1 & 66.67 & 66.67 & 77.78 \\
\bottomrule
\end{tabular}
\end{table*}

% =====================================================================
\subsection{False Positive Analysis}
\label{sec:eval:fp}

We analyse every false positive raised across the four categories and attribute
it to a root cause: SDK or marketing signals misread as a live surface, soft
404 / catch-all responses, or ambiguous backend-generated text.
Table~\ref{tab:rq1} and Table~\ref{tab:detection} together show that the
evidence gate and differential protocol confine false positives to
\textsc{XXX} cases, and we discuss the residual failures and their implications
for the evidence hierarchy here.
